% Appendix: PaTAS-TP theorems and proofs

\chapter{\ac{PaTAS-TP} Theorems and Proofs (\cref{sec:ptastheory})}\label{app:proof}

Throughout this appendix, opinions are binomial, $\omega = (b, d, u)$ with $b + d + u = 1$ and a common base rate $a = 1/2$ unless stated otherwise, and $q = (b_q, d_q, u_q)$ denotes a fixed evidence opinion. The revision sequences studied are of the form $\omega_{n+1} = \omega_n \circleddash q$ for a fusion operator $\circleddash$; the operators are those of \cref{sec:fusion}. None of them makes the set of opinions a group (\cref{sec:ptas_convergence}), so convergence is established operator by operator.

\section{Revision with a Constant Evidence Opinion}

\begin{theorem}[Averaging and weighted revision]\label{thm:abf}
Let $\omega_{n+1} = \omega_n \,\underline{\oplus}\, q$ (averaging belief fusion, \cref{def:averaging_fuse}) with $\omega_0$ arbitrary and $0 < u_0 \le 1$.
\begin{enumerate}
    \item If $q$ is not vacuous ($u_q < 1$), then $\omega_n \to q$, and the convergence is geometric: $1/u_{n} - 1/u_q = 2^{-n}\,(1/u_0 - 1/u_q)$ and $|b_{n+1} - b_q| = \frac{u_q}{u_n + u_q}\,|b_n - b_q|$, the contraction factor tending to $1/2$.
    \item If $q$ is vacuous ($u_q = 1$), then $\omega_n \to (0, 0, 1)$: the evidence of $\omega_0$ is halved at every step, $1/u_n - 1 = 2^{-n}\,(1/u_0 - 1)$.
\end{enumerate}
The same limit $q$ is obtained under weighted belief fusion $\hat{\oplus}$ (\cref{def:weighted_fuse}) when $q$ is not vacuous, whereas a vacuous $q$ leaves $\omega_0$ unchanged under weighted fusion.
\end{theorem}
\begin{proof}
For averaging fusion with $u_n + u_q > 0$,
\[
u_{n+1} = \frac{2\, u_n u_q}{u_n + u_q}, \qquad
b_{n+1} = \frac{b_n u_q + b_q u_n}{u_n + u_q}, \qquad
d_{n+1} = \frac{d_n u_q + d_q u_n}{u_n + u_q}.
\]
Taking reciprocals in the first relation gives $1/u_{n+1} = \tfrac12\,(1/u_n + 1/u_q)$, an affine recursion in $1/u_n$ with fixed point $1/u_q$ and contraction factor $1/2$, which proves the statements on $u_n$ in both cases. For the belief, $b_{n+1} - b_q = (b_n - b_q)\, u_q/(u_n + u_q)$, and $u_q/(u_n + u_q) \to 1/2 < 1$ when $u_q < 1$, so $b_n \to b_q$ geometrically; the same computation applies to $d_n$. When $u_q = 1$, $b_{n+1} = b_n/(u_n + 1)$ with $u_n \to 1$, so $b_n \to 0$, and likewise $d_n \to 0$. Note that a strictly decreasing distance to $q$ alone would not prove convergence \emph{to} $q$; the explicit contraction factors do. For weighted fusion the uncertainty obeys $u_{n+1} = (2 - u_n - u_q)\, u_n u_q / (u_n + u_q - 2 u_n u_q)$. For $u_q < 1$ this map is increasing on $(0,1]$ and its value lies strictly between $u_n$ and $u_q$, so $u_n \to u_q$ monotonically; the belief and disbelief recursions are convex combinations of $(b_n, d_n)$ and $(b_q, d_q)$ whose weight on $(b_n, d_n)$ tends to $1/2$, hence they converge to $(b_q, d_q)$ as above. When $u_q = 1$ the vacuous opinion is the neutral element of weighted fusion, so $\omega_n = \omega_0$.
\end{proof}

\begin{theorem}[Cumulative revision]\label{thm:cbf}
Let $\omega_{n+1} = \omega_n \oplus q$ (cumulative belief fusion, \cref{def:cumulative_fuse}) with $0 < u_0 \le 1$.
\begin{enumerate}
    \item If $q$ is not vacuous, then $1/u_n = 1/u_0 + n\,(1/u_q - 1)$, so $u_n \to 0$ like $1/n$, and $\omega_n$ converges to the dogmatic opinion $(b_q, d_q, 0)/(b_q + d_q)$: the initial opinion $\omega_0$ is forgotten.
    \item If $q$ is vacuous, then $\omega_n = \omega_0$ for all $n$, the vacuous opinion being the neutral element of cumulative fusion.
\end{enumerate}
\end{theorem}
\begin{proof}
Cumulative fusion adds evidence: in terms of the evidence parameters $r = W b/u$ and $s = W d/u$ of an opinion (Beta--opinion mapping with prior weight $W$), $\omega_n \oplus q$ has evidence $(r_n + r_q,\; s_n + s_q)$. Hence $r_n = r_0 + n\, r_q$, $s_n = s_0 + n\, s_q$ and $u_n = W/(W + r_n + s_n)$. Since $1/u = 1 + (r + s)/W$ and $r_q + s_q = W\,(1 - u_q)/u_q$, this gives the stated expression for $1/u_n$, and $u_n \to 0$ unless $r_q = s_q = 0$, that is, unless $q$ is vacuous. Then $b_n = r_n/(W + r_n + s_n) \to r_q/(r_q + s_q) = b_q/(b_q + d_q)$, and likewise for $d_n$. If $q$ is vacuous the evidence is unchanged. For example, $q = (0.5, 0, 0.5)$ drives every non-dogmatic $\omega_0$ to $(1, 0, 0)$.
\end{proof}

\begin{theorem}[Belief constraint revision]\label{thm:bcf}
Let $\omega_{n+1} = \omega_n \odot q$ (belief constraint fusion, \cref{def:constraint_fuse}) with $\omega_0$ non-dogmatic and $q$ neither vacuous nor totally conflicting with $\omega_0$. Then $\omega_n$ converges to $(1, 0, 0)$ if $b_q > d_q$, to $(0, 1, 0)$ if $d_q > b_q$, and, if $b_q = d_q$, to the dogmatic opinion $(\beta, 1 - \beta, 0)$ with $\beta = (b_0 + u_0)/(b_0 + d_0 + 2 u_0)$, fixed by the plausibilities of $\omega_0$. A vacuous $q$ leaves $\omega_0$ unchanged.
\end{theorem}
\begin{proof}
Belief constraint fusion is Dempster's rule restricted to binomial opinions. Write $P_n = b_n + u_n$ and $\bar P_n = d_n + u_n$ for the plausibilities of $x$ and $\bar x$ under $\omega_n$. The unnormalised fused masses are $b_n b_q + b_n u_q + u_n b_q$, $d_n d_q + d_n u_q + u_n d_q$ and $u_n u_q$, so the plausibilities multiply: $P_{n+1} \propto P_n P_q$, $\bar P_{n+1} \propto \bar P_n \bar P_q$ and $u_{n+1} \propto u_n u_q$, all with the same normalisation constant $1/(1 - K_n)$, $K_n = b_n d_q + d_n b_q$. Therefore
\[
\frac{P_n}{\bar P_n} = \frac{P_0}{\bar P_0}\Big(\frac{b_q + u_q}{d_q + u_q}\Big)^{n}, \qquad
\frac{u_n}{P_n} = \frac{u_0}{P_0}\Big(\frac{u_q}{b_q + u_q}\Big)^{n}.
\]
If $b_q > d_q$, the first ratio diverges, so $\bar P_n = d_n + u_n \to 0$; the second ratio tends to zero because $b_q > 0$, so $u_n \to 0$, hence $d_n \to 0$ and $b_n \to 1$. The case $d_q > b_q$ is symmetric. If $b_q = d_q > 0$, the first ratio is constant, equal to $P_0/\bar P_0$, while $u_n \to 0$, so the limit is dogmatic with $b/d = P_0/\bar P_0$, which gives $\beta$. A vacuous $q$ has $K = 0$ and returns $\omega_n$ unchanged.
\end{proof}

\begin{remark}[Amplification and numerical caveat]\label{rem:bcf-numerics}
The evidence $q = (0.4, 0.35, 0.25)$ drives the sequence to $(1, 0, 0)$: a slight majority of belief over disbelief is amplified to certainty. Iterating belief constraint fusion naively in floating point is moreover unstable, because each step multiplies the rounding error of $b + d + u$ by $1/(1 - K_n)$; in a test run with this $q$ the mass sum had drifted away from $1$ after $60$ steps and collapsed below $0.01$ after $100$ steps. An implementation that iterates this operator must renormalise after every step, as the implementation of this thesis does after every operator application.
\end{remark}

\section{The Implemented Update under Trusted Data}

\begin{theorem}[Implemented update under trusted data]\label{thm:implemented}
Let \cref{alg:trustupdate} run with averaging revision, base rate $a = 1/2$, $T_{lr} = (1, 0, 0)$, and a fully trusted conditioning factor, $T_x \oslash T_{y_\text{batch}} = (1, 0, 0)$, as holds for the first layer under trusted input and label opinions. Let $q_n = (b_{q_n}, d_{q_n}, u_q)$ be the node opinion of batch $n$; its uncertainty $u_q = W / (|\mathcal{N}(i)| + W)$ does not depend on $n$. Then, for any initial opinion with $u^{(0)} > 0$,
\begin{enumerate}
    \item $u^{(n)} \to u^* = u_q / 15$ monotonically, whatever the sequence $(q_n)$;
    \item $d^{(n+1)} = \lambda_n\, d^{(n)} + (1 - \lambda_n)\, d_{q_n}$ with $\lambda_n = u_q / (u^{(n)} + u_q) \to 15/16$, so that $d^{(n)}$ is a convex combination of the past batch evidence with geometrically decaying weights, and $b^{(n)} = 1 - d^{(n)} - u^{(n)}$;
    \item if $d_{q_n} \to d_q$, then $T_\theta^{(n)} \to (1 - d_q - u_q/15,\; d_q,\; u_q/15)$.
\end{enumerate}
\end{theorem}
\begin{proof}
With base rate $1/2$, binomial multiplication (\cref{def:binomial_multiplication}) by the dogmatic opinion $(1,0,0)$ maps $(b, d, u)$ to $(b + u/3,\; d,\; 2u/3)$. The three sub-steps of one batch act on the uncertainty as follows. The revision gives $u' = 2\, u\, u_q / (u + u_q)$. The learning-rate adjustment fuses $(b, d, u')$ with $(b + u'/3,\; d,\; 2u'/3)$ by averaging fusion, which gives $u'' = 2\, u' (2u'/3) / (u' + 2u'/3) = 4u'/5$. The auxiliary adjustment multiplies by $(1,0,0)$: $u''' = 2u''/3 = 8u'/15$. Composing the three maps, $u^{(n+1)} = \varphi(u^{(n)})$ with
\[
\varphi(u) = \frac{16}{15}\cdot\frac{u\, u_q}{u + u_q},
\]
which is increasing and concave on $(0, 1]$, satisfies $\varphi(u) > u$ exactly when $u < u_q/15$ and $\varphi(u) < u$ when $u > u_q/15$. Hence $u^{(n)}$ converges monotonically to $u^* = u_q/15$ from any positive start, and $\varphi'(u^*) = 15/16 < 1$ makes the convergence geometric. Only the revision changes the disbelief: multiplication by $(1,0,0)$ leaves $d$ unchanged, and the averaging fusion of two opinions with the same disbelief returns that disbelief. Thus $d^{(n+1)} = (d^{(n)} u_q + d_{q_n} u^{(n)}) / (u^{(n)} + u_q)$, which is the stated convex combination, with $\lambda_n \to u_q/(u^* + u_q) = 15/16$. The belief follows by normalisation, and the limit in the third statement is obtained by letting $n \to \infty$ in the two recursions.
\end{proof}

\begin{remark}[Worked example and the vacuous case]\label{rem:worked}
For a neuron with $30$ incoming edges and $W = 2$, $u_q = 1/16$ and $u^* = 1/240 \approx 0.0042$. With constant evidence of $r = 27$ weak and $s = 3$ strong gradients, $q = (0.844, 0.094, 0.063)$, and a direct iteration of the released implementation converges to $(0.9021, 0.0938, 0.0042)$ within about $60$ batches, as predicted. Under cumulative revision the same iteration reaches a numerically zero uncertainty within $20$ batches (the implementation does not clip small uncertainties; in single precision the value settles at the rounding floor of the normalisation $b + d + u = 1$, about $10^{-8}$) and then freezes, since an effectively dogmatic opinion is no longer changed by a less certain one; under belief constraint revision it reaches exactly $(1, 0, 0)$ within five batches. Under a vacuous label opinion the auxiliary sub-step multiplies by $(0, 0, 1)$, which maps $(b, d, u)$ to $(b/3,\; d,\; 1 - b/3 - d)$ and expands the uncertainty instead of contracting it; the same iteration then converges to $(0.147, 0.047, 0.806)$ under averaging revision and to $(0.210, 0.183, 0.607)$ under cumulative revision. In this regime the parameter uncertainty is large and the choice of the revision operator matters. For layers beyond the first, the input opinion $T_x$ entering the auxiliary sub-step is the propagated opinion of the previous layer, which is close to but not exactly dogmatic, so the contraction is weaker and the uncertainty settles above $u_q/15$: in the breast-cancer network of Experiment~1 the mean parameter uncertainty after training is $0.005$ in the first layer and $0.06$ in the output layer.
\end{remark}

\section{Multiplicative Revision}

\begin{theorem}[Convergence of the multiplicative sequence]\label{thm:geo}
Let $\omega_{n+1} = \omega_n \odot q$ with the binomial multiplication of \cref{def:binomial_multiplication} evaluated at a fixed base rate $a \in [0,1)$ for both arguments, and write $c = a/(1+a)$ ($c = 1/3$ for $a = 1/2$). Then:
\begin{itemize}
    \item $d_{n+1} = d_n + d_q - d_n d_q$, so $d_n \to 1$ if $d_q > 0$, and $d_n = d_0$ for all $n$ if $d_q = 0$;
    \item if $d_q > 0$, then $b_n \to 0$ and $u_n \to 0$, that is, $\omega_n \to (0, 1, 0)$;
    \item if $d_q = 0$, then $b_n \to b^* = c\,(1 - d_0)\, b_q \big/ \big((1 - c)(1 - b_q) + c\, b_q\big)$, which equals $(1 - d_0)\, b_q/(2 - b_q)$ for $a = 1/2$, and $u_n \to 1 - b^* - d_0$.
\end{itemize}
\end{theorem}
\begin{proof}
With equal base rates $a$ the multiplication reads $b_{n+1} = b_n b_q + c\,(b_n u_q + u_n b_q)$, $d_{n+1} = d_n + d_q - d_n d_q$ and $u_{n+1} = 1 - b_{n+1} - d_{n+1}$. The disbelief recursion is affine with fixed point $1$ and factor $1 - d_q$; if $d_q > 0$ it converges to $1$ monotonically, and if $d_q = 0$ it is constant. If $d_q > 0$, then $0 \le b_n \le 1 - d_n \to 0$ and $0 \le u_n \le 1 - d_n \to 0$, which proves the second statement without any approximation. If $d_q = 0$, then $u_q = 1 - b_q$, $d_n = d_0$ and $u_n = 1 - b_n - d_0$, so the belief recursion becomes the affine map $b_{n+1} = \rho\, b_n + c\,(1 - d_0)\, b_q$ with $\rho = (1 - c)\, b_q + c\,(1 - b_q)$. Since $0 \le c < 1/2$, $\rho$ is a convex combination of $b_q$ and $1 - b_q$ and satisfies $\rho < 1$, so $b_n$ converges to $b^* = c\,(1 - d_0)\, b_q / (1 - \rho)$, which is the stated value; $u_n$ follows by normalisation.
\end{proof}

\section{Proof of \cref{thm:convergence}}
\begin{proof}
Statement (i) collects \cref{thm:abf,thm:cbf,thm:bcf}. Statement (ii) is \cref{thm:implemented}: under trusted data the node opinion $q_n$ enters the update only through the revision, and the theorem shows that the composed map contracts the uncertainty to $u_q/15$ independently of $(q_n)$ and turns the disbelief into a convex combination of the past $d_{q_n}$ with weights $(1 - \lambda_n)\prod_{k > n}\lambda_k$, which decay geometrically once $\lambda_n$ is close to $15/16$. If $d_{q_n}$ converges, so does this running mean, which gives the limit; if $d_{q_n}$ is stationary but fluctuating, the running mean stays within the range of the recent evidence and inherits its fluctuations, damped by the window. Statement (iii) is \cref{thm:geo}. The assumptions of stable input and label opinions enter through the auxiliary sub-step: they fix the factor $T_x \oslash T_{y_\text{batch}}$, so that the same map is applied at every batch and the fixed-point argument applies.
\end{proof}

\section{Proof of \cref{thm:infvac}}

Let \( \omega^\emptyset = (0, 0, 1, a) \) denote a vacuous binomial opinion over any variable, with arbitrary base rate \( a \in [0, 1] \). Then the following two properties hold:

\begin{enumerate}
    \item \emph{Discounting a Vacuous Opinion Yields a Vacuous Opinion.} \\
    For any trust value binomial opinion \(\omega^A_B \), the discounted opinion
    \[
    \omega^{[A;B]}_X = \omega^A_B \otimes (0,0,1,a)= (0,0,1,a)
    \]
    \textit{Proof.} Using the trust discounting operator from Subjective Logic, we set $P$ to the projected probability of \(\omega^A_B\):
    \[\omega^{[A;B]}_X =
    \begin{cases}
    b^{[A;B]}_X = P\cdot 0 = 0, \\
    d^{[A;B]}_X = P\cdot 0 = 0, \\
    u^{[A;B]}_X = 1 - b^{[A;B]}_X-d^{[A;B]}_X = 1, \\
    a^{[A;B]}_X = a.
    \end{cases}
    \]
    Hence, \( \omega^{[A;B]} = (0,0,1,a) \).

    \item \emph{PaTAS Feedforward on Vacuous Input Yields Vacuous Output.} \\
    Let \( T_x = (0,0,1,a)\) be the trust assessment of an input to PaTAS. Then for any parameter trust configuration \( T_\theta \), the output trust assessment satisfies:
    \[
    T_y = \text{IPTA}(T_x) = (0,0,1,a).
    \]
    \textit{Proof.} The PaTAS feedforward mechanism computes for each neuron:
    \[
    T_z = \bigoplus_i \left( T_{x_i} \otimes T_{\theta_i} \right).
    \]
    Since each \( T_{x_i} = (0,0,1,a) \), and using the result from part (1), we get:
    \[
    T_{x_i} \otimes T_{\theta_i} = (0,0,1,a), \quad \forall i.
    \]
    Then, by fusion of vacuous opinions (averaging fusion is idempotent, so the mean of vacuous opinions is vacuous):
    \[
    T_z = \bigoplus_i (0,0,1,a) = (0,0,1,a).
    \]
    This holds recursively through all layers of PaTAS, including the output layer, hence:
    \[
    T_y = (0,0,1,a).
    \]
\end{enumerate}

\section{Proof of \cref{thm:sym}}
\begin{proof}
We prove the theorem in two steps.

\emph{(1) Symmetry of the Discount Operator:} \\
Let $\omega_\theta = (b_\theta, d_\theta, u_\theta)$ be a binomial opinion representing the trust weight associated with a connection in the PaTAS. Let $P$ denote the projected probability of $\omega_\theta$, defined as:
\[
P = b_\theta + a u_\theta,
\]
where $a$ is the base rate (typically $a = 0.5$ for binary domains).

Let $x = (b, d, u)$ be any binomial opinion and $\bar{x} = (d, b, u)$ its symmetric counterpart. Then the trust discounting operation yields:
\[
\omega_\theta \otimes x = (P \cdot b, P \cdot d, 1 - P \cdot (b + d)),
\]
\[
\omega_\theta \otimes \bar{x} = (P \cdot d, P \cdot b, 1 - P \cdot (b + d)).
\]
Since $b + d = 1 - u$, both discounted opinions have identical uncertainty and symmetric belief/disbelief masses. Hence, $\omega_\theta \otimes x$ and $\omega_\theta \otimes \bar{x}$ are symmetric.

\emph{(2) Symmetry Preservation under Fusion:} \\
Let $\mathcal{X} = \{x_1, \dots, x_n\}$ be a set of discounted opinions resulting from symmetric inputs, and let $\bar{\mathcal{X}} = \{\bar{x}_1, \dots, \bar{x}_n\}$ be their symmetric counterparts. Consider any fusion operator $\bigoplus$ that treats belief and disbelief symmetrically (the component-wise mean used in this thesis, as well as averaging, cumulative or consensus fusion in Subjective Logic) applied to $\mathcal{X}$ and $\bar{\mathcal{X}}$.

Since each pair $(x_i, \bar{x}_i)$ is symmetric and the operator treats belief and disbelief symmetrically, we have:
\[
\bigoplus_{i=1}^n x_i = (b', d', u') \quad \Rightarrow \quad \bigoplus_{i=1}^n \bar{x}_i = (d', b', u').
\]
Thus, the output of the PaTAS feedforward inference remains symmetric when symmetric inputs are provided.

\emph{(3) Invariance under Fusion:} \\
If all input opinions assign the value zero to a specific mass (belief or disbelief), the fused result preserves that value under every operator considered here, since the fused belief (disbelief) is a weighted average of the input beliefs (disbeliefs). In particular, if all input opinions have belief mass equal to zero, the fused result also has belief mass zero. The same applies to disbelief mass.

Therefore, if the discounting step yields discounted opinions with zero belief (or zero disbelief) across all components, then the fusion stage will preserve that zero mass in the final output.

\emph{(4) Fully Trusted and Distrusted Cases:} \\
For $x = (1, 0, 0)$ (fully trusted), we have:
\[
\omega_\theta \otimes x = (P, 0, 1 - P),
\]
so the discounted opinion assigns zero disbelief. As noted above, the fusion of such discounted opinions will also assign zero disbelief.

For $\bar{x} = (0, 1, 0)$ (fully distrusted), we have:
\[
\omega_\theta \otimes \bar{x} = (0, P, 1 - P),
\]
so the discounted opinion assigns zero belief. Therefore, the belief from a fully distrusted input is always zero in the PaTAS output.
\end{proof}
